\documentclass[a4paper]{article}

\usepackage[spanish]{babel}
\usepackage{graphicx}
\usepackage{amsmath, amssymb}
\usepackage[margin=2cm]{geometry}
\usepackage{fancyhdr}
\usepackage{enumerate}
\usepackage[shortlabels]{enumitem}
\usepackage{parskip}
\usepackage[most]{tcolorbox}
\usepackage[hidelinks]{hyperref}
\usepackage{float}
\usepackage{booktabs}
\usepackage{mathrsfs}


% cabecera
\pagestyle{fancy}
\fancyhead[l]{Física Matemática}
\fancyhead[c]{Parcial 1}
\fancyhead[r]{\today}
\fancyfoot[c]{\thepage}
\renewcommand{\headrulewidth}{0.2pt} % linea horizontal


\begin{document}


\section{Problema 1: potencial eléctrico}
En coordenadas polares $(r'=ra_0,\theta)$ el potencial eléctrico de un electrón en el nivel $1s$ en un átomo de hidrógeno está dado por la siguiente expresión

$$V(r) = \frac{1}{4\pi\varepsilon_0}\frac{q}{a_0}\left(\frac{1}{2r} \gamma(3,2r) + \Gamma(2,2r)\right)$$ 

donde $r$ es un número en términos del radio de Bohr $a_0$ y $\gamma(x, y)$, $\Gamma(x, y)$ son las funciones gamma incompletas. Determine el potencial eléctrico para $r<<1$ a orden $\mathcal{O}(r^2)$.

\textit{Ayuda:}

$$\gamma(n, x) = \int_{0}^{x}e^{-t}t^{n-1}dt = (n-1)!\left(1-e^{-x}\sum_{s=0}^{n-1} \frac{x^s}{s!}\right)$$

$$\Gamma(n, x) = \int_{x}^{\infty}e^{-t}t^{n-1}dt = (n-1)!e^{-x}\sum_{s=0}^{n-1} \frac{x^s}{s!}$$

\section{Problema 2: Oscilaciones}

Una partícula de masa $m$ en un oscilador tiene una energía potencial $V(x)=A|x|^n$

\begin{enumerate}[a.]
	\item Mostrar que el período del oscilador es $$\tau = 2\sqrt{2m}\int_{0}^{x_{\text{max}}}\frac{dx}{\sqrt{E-Ax^n}}$$ donde $x_{\text{max}}$ es el punto clásico de retorno (donde la velocidad es cero) y $A|x_{\text{max}}|^n=E$ es la energía mecánica del sistema.
	\item Determinar $\tau$. \textit{Ayuda:} $$B(m+1,n+1) = 2\int_{0}^{\pi/2} \cos^{2m+1}\theta \sin^{2n+1}\theta d\theta=\frac{m!n!}{(m+n+1)!}$$
	\item Verificar que para el caso de un movimiento armónico simple (\textit{M.A.S}), donde $V(x) = \left(\frac{k}{2}\right)x^2$ entonces $\tau = 2\pi \sqrt{\frac{m}{k}}$
\end{enumerate}

\section{Problema 3: Expandir función con base de Fourier discreta}
Expandir la función 

$$
f(x) =
\begin{cases}
-x, & \text{si }  -\pi \leq x \leq 0 \\
x, & \text{si } 0 < x \leq \pi
\end{cases}
$$

en la base discreta, ortonormal de Fourier $\{\phi_n(x)\} = 	\left\{\frac{1}{\sqrt{2L}}e^{i n\pi x /L}\right\}$ en $c\leq x \leq c+2L$, $n=-\infty, .., 0, ..., \infty$ del espacio de Hilbert de funciones de cuadrado integrable $\mathscr{L}^2$ con peso $p(x)=1$


\section{Problema 4: Análisis de Fourier para un pulso de sonido}

Considere un pulso de sonido de ancho $2T$ y amplitud constante de $1$ dado por

$$
f(t) =
\begin{cases}
0, & \text{si }  T \leq x \leq -T \\
1, & \text{si } -T < x \leq T
\end{cases}
$$

\begin{enumerate}[a.]
	\item Usar el análisis de Fourier con la base ortogonal $\{\phi(\omega, t)\} = \left\{\frac{1}{\sqrt{2\pi}}e^{i \omega t}\right\}$ de $\mathscr{L}^2$ para estimar el ancho de banda $\Delta \omega$ mínimo. Analizar el caso $\Delta t = 2T=1\text{s}$
	\item ¿Cuál es el contenido de frecuencias en el pulso? 
	
\end{enumerate}


\bibliographystyle{plainurl}
\bibliography{bibliografia} 
\end{document}
